% !TEX program = pdflatex
% Modelo: Apresentacao 01102026 - Hipoteses Modelos Dados.
% Fontes: aed_chap1_pt_ANALISE.ipynb e dados_agregados_aed_chap1_pt_ANALISE.md.
% Rastreabilidade, compilacao e revisao: 2_INTI/audit/apresentacao_15102026/.
\documentclass[aspectratio=169,11pt]{beamer}
\usetheme{Madrid}
\usepackage[T1]{fontenc}
\usepackage[utf8]{inputenc}
\usepackage[brazilian]{babel}
\usepackage{lmodern}
\usepackage{booktabs,array,tabularx}
\usepackage{amsmath}
\usepackage{pgfplots}
\usepgfplotslibrary{statistics}
\pgfplotsset{compat=1.18}
\definecolor{verde}{HTML}{24574D}
\definecolor{azul}{HTML}{1769A5}
\definecolor{grafite}{HTML}{253A40}
\definecolor{laranja}{HTML}{B87532}
\setbeamercolor{structure}{fg=verde}
\setbeamercolor{title}{bg=verde,fg=white}
\setbeamercolor{frametitle}{bg=verde,fg=white}
\setbeamercolor{normal text}{fg=grafite}
\setbeamertemplate{navigation symbols}{}
\setbeamertemplate{footline}{\leavevmode\hbox{\begin{beamercolorbox}[wd=.86\paperwidth,ht=2.7ex,dp=1.2ex,leftskip=1em]{author in head/foot}Retornos, instituições e conversão da vegetação nativa\end{beamercolorbox}\begin{beamercolorbox}[wd=.14\paperwidth,ht=2.7ex,dp=1.2ex,center]{date in head/foot}\insertframenumber\,/\,\inserttotalframenumber\end{beamercolorbox}}}
\setbeamersize{text margin left=8mm,text margin right=8mm}
\setbeamertemplate{itemize items}[circle]
\newcommand{\fonte}[1]{\par\vfill{\fontsize{7.5}{9}\selectfont\color{gray}#1\par}}
\newcommand{\tpar}{\par\vspace{2mm}}
\newcolumntype{P}[1]{>{\raggedright\arraybackslash}p{#1}}
\newcolumntype{L}{>{\raggedright\arraybackslash}X}
\renewcommand{\arraystretch}{1.16}
\title[Retornos, instituições e conversão]{Retornos Econômicos, Instituições e Conversão da Vegetação Nativa em Mato Grosso}
\subtitle{Dados preliminares — 2001--2024}
\author{Fernando Souza da Vieira}
\institute{Tópicos em Microeconomia}
\date{15 de outubro de 2026}
\begin{document}

% 1 — Capa; título e autoria de chap1/main.tex.
\begin{frame}
  \titlepage
\end{frame}

% 2 — Notebook, célula 05ad2c49.
\begin{frame}{Problema de pesquisa}
\begin{block}{Pergunta central}
Como os retornos econômicos esperados e o ambiente institucional condicionam a taxa de conversão da vegetação nativa em Mato Grosso, e em que medida essa relação varia conforme a capacidade produtiva local da terra, as atividades agropecuárias e os biomas?
\end{block}
\medskip
\textbf{Objeto da análise}\par\smallskip
A conversão direta de vegetação nativa para agricultura e pastagem, observada por município e ano.
\end{frame}

% 3 — Notebook, células a4505b2a e 4fda2f9d.
\begin{frame}{Hipótese e função}
\begin{block}{H1 — Retorno esperado e conversão}
A conversão de vegetação nativa para usos agropecuários aumenta quando se eleva o retorno econômico esperado desses usos.
\end{block}
\medskip
\begin{equation*}
\begin{aligned}
\text{Conversão da vegetação nativa}
  = f\bigl(&\text{preços agropecuários},\\
           &\text{produtividade da terra},\\
           &\text{ambiente institucional}\bigr)
\end{aligned}
\end{equation*}
\end{frame}

% 4 — Notebook, célula 0d071678.
\begin{frame}{O que o artigo busca explicar}
O artigo busca explicar, compreender e mensurar se:
\medskip
\begin{enumerate}
\setlength{\itemsep}{4mm}
\item Aumentos nos preços dos produtos agropecuários se traduzem em maior conversão onde a \textbf{capacidade produtiva da terra} para esses usos é mais elevada.
\item As \textbf{condições institucionais} modificam a intensidade dessa resposta.
\item Essas relações entre expectativas e respostas aos ambientes econômico e institucional variam conforme \textbf{atividades agropecuárias e biomas}.
\end{enumerate}
\fonte{Nesta etapa, a AED descreve o fenômeno e a cobertura dos dados; ainda não testa H1.}
\end{frame}

% 5 — Notebook, df1.shape/info e descrição quantitativa/qualitativa.
\begin{frame}{df1: visão geral do painel completo}
\begin{block}{Dimensão e unidade de observação}
\textbf{3.384 linhas $\times$ 296 colunas}: 141 municípios de Mato Grosso, acompanhados por 24 anos (2001--2024). Cada linha representa um município--ano.
\end{block}
\small
\begin{tabularx}{\textwidth}{P{3.15cm} L}
\toprule
\textbf{Componente} & \textbf{Conteúdo quantitativo e qualitativo}\\
\midrule
Medidas numéricas & Conversão, estoques, produção, capacidade, preços, registros institucionais, chuva e relevo.\\
Categorias e identificação & Código municipal, ano, malha territorial e identificação das fontes.\\
Controle de qualidade & Indicadores de cobertura, parcialidade e pendências por fonte.\\
\bottomrule
\end{tabularx}\par
\smallskip
\textbf{Cobertura:} 2.772 valores válidos de Y e 612 ausências; chuva anual com 3.384 valores.
\end{frame}

% 6 — Notebook, agregado-1a-000 a 032.
\begin{frame}{df1a: perfis municipais agregados de 2001--2024}
\begin{block}{Agregação temporal}
\textbf{141 linhas $\times$ 105 colunas}. Cada linha resume um município; as estatísticas entre municípios atribuem o mesmo peso a cada um.
\end{block}
\small
\begin{tabularx}{\textwidth}{P{3.25cm} L}
\toprule
\textbf{Tipo de informação} & \textbf{Regra e conteúdo do agregado}\\
\midrule
Conversão & Somas observadas e totais completos; taxas médias e ponderadas; número de anos válidos.\\
Medidas anuais & Médias e dispersão nos anos disponíveis, com cobertura própria de cada medida.\\
Características fixas & Capacidade histórica PAM, frações de bioma, elevação e oito medidas/cenários GAEZ.\\
Estoques datados & Vegetação, agricultura e pastagem em 2001 e 2024; estoques nativos do par de mapas em 2000 e 2023.\\
\bottomrule
\end{tabularx}\par
\smallskip
Ausências são preservadas. O total de conversão em 2001--2024 exige \textbf{24 anos válidos}.
\end{frame}

% 7 — Relatório, seção 3; agregado-1a-048.
\begin{frame}{df1a: estoques mapeados de uso da terra}
\begin{columns}[T,onlytextwidth]
\column{.53\textwidth}
% Gerado de CSVs agregados; dados numericos incorporados para portabilidade.
% Notebook SHA-256: e8ce396c21be212f5a12d673622adf543234983c3cd142088ea2dbe8e7621004
% Parecer SHA-256: 1ee24339005b98f24afb2375d080f1aaa8ff1ab0e8812427b8b103b890cef98a
% Requer: pgfplots, biblioteca statistics (somente taxas.tex).
% Ajuste opcional no chamador: \def\AlturaGrafico{4.4cm} e \def\LarguraGrafico{\linewidth}.
\begingroup
\providecommand{\AlturaGrafico}{4.4cm}
\providecommand{\LarguraGrafico}{\linewidth}
\definecolor{grafverde}{HTML}{24574D}
\definecolor{grafazul}{HTML}{1769A5}
\definecolor{grafgrafite}{HTML}{253A40}
\definecolor{graflaranja}{HTML}{B97736}
\begin{tikzpicture}
\begin{axis}[width=\LarguraGrafico,height=\AlturaGrafico,
font=\scriptsize,tick label style={font=\scriptsize},
label style={font=\scriptsize},
axis line style={grafgrafite!65},tick style={grafgrafite!65},
grid style={black!10},
legend style={font=\scriptsize,draw=none,fill=none,cells={anchor=west}},
/pgf/number format/use comma,
ybar,bar width=11pt,ymin=0,ymax=75,ytick={0,15,30,45,60,75},
symbolic x coords={Nativa,Agricultura,Pastagem},xtick=data,
enlarge x limits=0.22,ymajorgrids=true,
ylabel={Milh\~oes de hectares},
legend columns=2,legend style={at={(0.5,-0.20)},anchor=north},
]
\addplot+[fill=grafverde,draw=grafverde,area legend] coordinates {(Nativa,64.6733010293) (Agricultura,5.0039899540) (Pastagem,18.3858676477)};
\addplot+[fill=grafazul,draw=grafazul,area legend] coordinates {(Nativa,54.9390637610) (Agricultura,12.4175621478) (Pastagem,20.2049180566)};
\legend{2001,2024}
\end{axis}
\end{tikzpicture}
\endgroup

\column{.44\textwidth}
\small
\textbf{Estoques em milhões de hectares}\par\smallskip
\setlength{\tabcolsep}{3pt}
\begin{tabularx}{\linewidth}{L r r r}
\toprule
 & \textbf{2001} & \textbf{2024} & \textbf{$\Delta$\%}\\
\midrule
Nativa & 64,673 & 54,939 & $-15,05$\\
Agricultura & 5,004 & 12,418 & $+148,15$\\
Pastagem & 18,386 & 20,205 & $+9,89$\\
\bottomrule
\end{tabularx}\par
\medskip
A agricultura cresce nos \textbf{141 municípios}. A pastagem aumenta em 75 e diminui em 66.
\end{columns}
\end{frame}

% 8 — Preços são série contextual separada do agregado; GAEZ é parte de df1a.
\begin{frame}{df1a: capacidade produtiva e contexto de preços}
\begin{columns}[T,onlytextwidth]
\column{.55\textwidth}
\small\textbf{Preços internacionais reais}\par\smallskip
{\def\LarguraGrafico{.96\linewidth}% Gerado de CSVs agregados; dados numericos incorporados para portabilidade.
% Notebook SHA-256: e8ce396c21be212f5a12d673622adf543234983c3cd142088ea2dbe8e7621004
% Parecer SHA-256: 1ee24339005b98f24afb2375d080f1aaa8ff1ab0e8812427b8b103b890cef98a
% Requer: pgfplots, biblioteca statistics (somente taxas.tex).
% Ajuste opcional no chamador: \def\AlturaGrafico{4.4cm} e \def\LarguraGrafico{\linewidth}.
\begingroup
\providecommand{\AlturaGrafico}{4.4cm}
\providecommand{\LarguraGrafico}{\linewidth}
\definecolor{grafverde}{HTML}{24574D}
\definecolor{grafazul}{HTML}{1769A5}
\definecolor{grafgrafite}{HTML}{253A40}
\definecolor{graflaranja}{HTML}{B97736}
\begin{tikzpicture}
\begin{axis}[width=\LarguraGrafico,height=\AlturaGrafico,
font=\scriptsize,tick label style={font=\scriptsize},
label style={font=\scriptsize},
axis line style={grafgrafite!65},tick style={grafgrafite!65},
grid style={black!10},
legend style={font=\scriptsize,draw=none,fill=none,cells={anchor=west}},
/pgf/number format/use comma,
xmin=2001,xmax=2024,ymin=80,ymax=280,
xtick={2001,2006,2012,2018,2024},ytick={100,150,200,250},
scaled x ticks=false,xticklabel style={/pgf/number format/1000 sep=},
ymajorgrids=true,ylabel={\'Indice real (2001 = 100)},
legend columns=3,legend style={at={(0.5,-0.20)},anchor=north},
]
\addplot[black!30,densely dotted,forget plot] coordinates {(2001,100) (2024,100)};
\addplot[grafverde,thick,solid,no marks] coordinates {(2001,100.0000000000) (2002,106.8290732597) (2003,129.7117249281) (2004,146.9069763879) (2005,127.2378249174) (2006,120.4844300888) (2007,166.8318467877) (2008,218.8594646600) (2009,178.4591575935) (2010,185.3568891076) (2011,216.1428785961) (2012,234.3743788098) (2013,214.0360991498) (2014,185.2254487202) (2015,149.5971115163) (2016,152.7121258100) (2017,145.1431452906) (2018,142.0389994820) (2019,130.5314953683) (2020,141.9844369982) (2021,194.7813080281) (2022,208.7129059518) (2023,177.4885912974) (2024,133.3578613041)};
\addlegendentry{Soja}
\addplot[grafazul,thick,dashed,no marks] coordinates {(2001,100.0000000000) (2002,108.9311091293) (2003,113.1334925479) (2004,116.9792709047) (2005,99.7806603170) (2006,119.3350954823) (2007,155.9036074226) (2008,204.3366638623) (2009,152.3856266710) (2010,168.2986028315) (2011,256.0544345099) (2012,256.6122596652) (2013,219.9505085100) (2014,160.9161540693) (2015,141.4483008255) (2016,130.9809176463) (2017,124.5230137193) (2018,129.3172314256) (2019,131.3674275253) (2020,126.2060401501) (2021,189.1878054158) (2022,215.0929562744) (2023,163.9505286597) (2024,120.0016960720)};
\addlegendentry{Milho}
\addplot[graflaranja,thick,dashdotted,no marks] coordinates {(2001,100.0000000000) (2002,98.6555269731) (2003,97.2413579263) (2004,118.9469679035) (2005,117.7310104897) (2006,110.5633897398) (2007,109.2941157898) (2008,126.1821158398) (2009,106.9681582852) (2010,131.8237306951) (2011,152.4882865556) (2012,154.0790058149) (2013,149.0635684869) (2014,179.3000325991) (2015,161.8834998043) (2016,142.2702304421) (2017,150.5115941347) (2018,140.5315195989) (2019,156.6631481360) (2020,151.7037993815) (2021,165.6229252533) (2022,161.7078089666) (2023,135.2951361158) (2024,158.9944278889)};
\addlegendentry{Carne}
\end{axis}
\end{tikzpicture}
\endgroup
}
\column{.42\textwidth}
\small
\textbf{GAEZ (Global Agro-Ecological Zones): rendimento potencial}\par
\footnotesize Média municipal, kg de matéria seca/ha\par\smallskip
\small
\begin{tabularx}{\linewidth}{L r r}
\toprule
\textbf{Cultura} & \textbf{Alto insumo} & \textbf{Baixo}\\
\midrule
Soja & 2.921,53 & 302,57\\
Milho & 3.303,27 & 890,83\\
\bottomrule
\end{tabularx}\par
\medskip
Os preços variam no tempo; o potencial depende da cultura e do cenário de produção.
\tpar
GAEZ tem 141 municípios; a capacidade histórica PAM (Pesquisa Agrícola Municipal) tem 47 para soja e 117 para milho.
\end{columns}
\fonte{Preços: série anual separada, dólares de 2024, índice 2001 = 100.}
\end{frame}

% 9 — Parecer, seções 3 a 6.
\begin{frame}{df1a: interpretação econômica dos agregados}
\begin{itemize}
\setlength{\itemsep}{4mm}
\item A expansão agrícola e as trajetórias distintas da pastagem indicam \textbf{reorganização dos usos da terra}. A substituição entre usos já instalados é uma possibilidade a investigar com transições específicas.
\item Os ciclos de preços podem alterar os incentivos de forma diferente conforme a atividade e a capacidade local. \textbf{Preço internacional não é retorno líquido esperado}: custos e transmissão ao produtor ainda precisam ser considerados.
\item PAM observada, capacidade histórica e cenários GAEZ descrevem \textbf{objetos distintos}. Maior rendimento físico não implica, por si só, maior rentabilidade.
\end{itemize}
\fonte{Os agregados descrevem transformações e condições produtivas, não identificam seus mecanismos causais.}
\end{frame}

% % 10 — Parecer, seções 5, 6 e 9.
% \begin{frame}{Como df1a se relaciona com o artigo}
% \small
% \begin{tabularx}{\textwidth}{P{3.65cm} L}
% \toprule
% \textbf{Evidência descritiva} & \textbf{Relação com a pergunta do artigo}\\
% \midrule
% Estoques e fluxos distintos & Separar expansão da atividade em áreas já utilizadas de abertura direta de vegetação nativa.\\[2mm]
% Capacidade produtiva heterogênea & Investigar se o mesmo movimento de preço se associa a respostas diferentes conforme as condições locais.\\[2mm]
% Preços comuns aos municípios & Recuperar a variação anual: a média de todo o período não preserva o movimento dos incentivos.\\[2mm]
% Medidas e coberturas diferentes & Definir uma medida de capacidade defensável e verificar quais municípios entram na comparação.\\
% \bottomrule
% \end{tabularx}\par
% \medskip
% \normalsize\textbf{df1a contextualiza os territórios; a avaliação de H1 depende do painel anual.}
% \fonte{Os resultados atuais não confirmam nem rejeitam a hipótese.}
% \end{frame}

% 11 — Notebook, df2-criterios, df2-preliminar, shape/info.
\begin{frame}{df2: visão geral do painel selecionado}
\begin{block}{Seleção para a análise preliminar de H1}
\textbf{3.384 linhas $\times$ 30 colunas}: os mesmos 141 municípios e 24 anos de df1, com 2 identificadores, 24 medidas e 4 indicadores de qualidade.
\end{block}
\small
\begin{tabularx}{\textwidth}{P{3.1cm} L}
\toprule
\textbf{Dimensão} & \textbf{Conteúdo selecionado}\\
\midrule
Conversão e produção & Conversão para agricultura e pastagem, vegetação inicial, rendimentos, rebanho e lotação aparente.\\
Incentivos e território & Preços reais, capacidade histórica PAM, frações de bioma, chuva e elevação.\\
Instituições e qualidade & Autos e multas reais de desmatamento; cobertura de Y, ausência de registro identificado e parcialidade dos registros.\\
\bottomrule
\end{tabularx}\par
\smallskip
\textbf{Limites:} GAEZ e custos não estão nessa seleção; as lacunas de informação permanecem.
\fonte{Os indicadores de qualidade não são medidas de efetividade institucional, apenas de ação.}
\end{frame}

% 12 — df2a, corrigindo a repetição de df1a no roteiro solicitado.
\begin{frame}{df2a: agregação e cobertura da conversão}
\begin{block}{Perfil municipal selecionado}
\textbf{141 linhas $\times$ 88 colunas}. As 88 colunas comuns coincidem com df1a: são as mesmas medidas e regras de cálculo, não uma confirmação independente.
\end{block}
\begin{columns}[T,onlytextwidth]
\column{.40\textwidth}
\small
\begin{tabularx}{\linewidth}{L r}
\toprule
\textbf{Cobertura de Y} & \textbf{Municípios}\\
\midrule
24 anos completos & 95\\
Cobertura parcial & 42\\
Sem Y válido & 4\\
\bottomrule
\end{tabularx}\par
\smallskip
Os completos representam 67,4\% dos municípios, sem garantia de representatividade territorial.
\column{.57\textwidth}
\small\textbf{Taxa anual média ponderada por município}
\[
100\,\frac{\sum_{t\in T_i}\text{conversão}_{it}}{\sum_{t\in T_i}\text{vegetação nativa inicial}_{it}}
\]
Numerador e denominador usam os mesmos anos válidos, por destino.
\tpar
A taxa resume a intensidade anual; não mede a perda acumulada da vegetação inicial.
\end{columns}
\end{frame}

% 13 — Notebook, agregado-2a-030/036, 95 municípios completos.
\begin{frame}{df2a: conversão direta por destino}
\begin{columns}[T,onlytextwidth]
\column{.49\textwidth}
% Gerado de CSVs agregados; dados numericos incorporados para portabilidade.
% Notebook SHA-256: e8ce396c21be212f5a12d673622adf543234983c3cd142088ea2dbe8e7621004
% Parecer SHA-256: 1ee24339005b98f24afb2375d080f1aaa8ff1ab0e8812427b8b103b890cef98a
% Requer: pgfplots, biblioteca statistics (somente taxas.tex).
% Ajuste opcional no chamador: \def\AlturaGrafico{4.4cm} e \def\LarguraGrafico{\linewidth}.
\begingroup
\providecommand{\AlturaGrafico}{4.4cm}
\providecommand{\LarguraGrafico}{\linewidth}
\definecolor{grafverde}{HTML}{24574D}
\definecolor{grafazul}{HTML}{1769A5}
\definecolor{grafgrafite}{HTML}{253A40}
\definecolor{graflaranja}{HTML}{B97736}
\begin{tikzpicture}
\begin{axis}[width=\LarguraGrafico,height=\AlturaGrafico,
font=\scriptsize,tick label style={font=\scriptsize},
label style={font=\scriptsize},
axis line style={grafgrafite!65},tick style={grafgrafite!65},
grid style={black!10},
legend style={font=\scriptsize,draw=none,fill=none,cells={anchor=west}},
/pgf/number format/use comma,
boxplot/draw direction=y,xmin=0.4,xmax=2.6,ymin=0,ymax=5,
xtick={1,2},xticklabels={Agricultura,Pastagem},
ytick={0,1,2,3,4,5},ymajorgrids=true,
ylabel={Taxa anual m\'edia ponderada (\%)},
legend style={at={(0.5,-0.20)},anchor=north},
]
\addplot+[
draw=grafverde,fill=grafverde!25,mark=*,mark size=1.5pt,
mark options={fill=grafverde,draw=grafverde},forget plot,
boxplot prepared={draw position=1,
lower whisker=0.000450151748,lower quartile=0.008456447345,
median=0.030110746995,upper quartile=0.134456746165,upper whisker=0.271390174812}
] coordinates {(0,0.324094880044) (0,0.351016934504) (0,0.374321853199) (0,0.433539648275) (0,0.470134352143) (0,0.473911493134) (0,0.514120356654) (0,0.516830134897) (0,0.626592442615) (0,0.998122141341) (0,1.528089212448)};
\addplot+[
draw=grafazul,fill=grafazul!25,mark=*,mark size=1.5pt,
mark options={fill=grafazul,draw=grafazul},forget plot,
boxplot prepared={draw position=2,
lower whisker=0.096504047886,lower quartile=0.623779749666,
median=1.079040156580,upper quartile=1.495796115814,upper whisker=2.622435998947}
] coordinates {(0,2.992642046469) (0,3.185403774944) (0,3.331745277705) (0,4.294436940632) (0,4.458760895268)};
\addplot[only marks,mark=diamond*,mark size=2.5pt,graflaranja] coordinates {(1,0.118496647617) (2,1.197874724151)};
\addlegendentry{M\'edia entre munic\'ipios}
\end{axis}
\end{tikzpicture}
\endgroup

\column{.48\textwidth}
\small\textbf{Taxas anuais médias ponderadas (\%)}\par\smallskip
\begin{tabularx}{\linewidth}{L r r}
\toprule
 & \textbf{Agricultura} & \textbf{Pastagem}\\
\midrule
Média & 0,1185 & 1,1979\\
Mediana & 0,0301 & 1,0790\\
1º quartil & 0,0085 & 0,6238\\
3º quartil & 0,1345 & 1,4958\\
\bottomrule
\end{tabularx}\par
\medskip
A taxa de pastagem supera a agrícola em \textbf{93 dos 95 municípios}.
\tpar
A média para pastagem é \textbf{10,1 vezes} a agrícola; a distribuição da agricultura é mais assimétrica.
\end{columns}
\fonte{95 municípios com 24 anos válidos; parecer, seção 2. Caixas: quartis e mediana; hastes até 1,5 intervalo interquartil; pontos além das hastes são observações, não erros identificados.}
\end{frame}

% 14 — Cobertura 2a e parecer, seções 1 e 7.
\begin{frame}{df2a: disponibilidade e registros institucionais}
\begin{columns}[T,onlytextwidth]
\column{.52\textwidth}
\small\textbf{Municípios com dados, por medida}\par\smallskip
{\def\LarguraGrafico{.94\linewidth}% Gerado de CSVs agregados; dados numericos incorporados para portabilidade.
% Notebook SHA-256: e8ce396c21be212f5a12d673622adf543234983c3cd142088ea2dbe8e7621004
% Parecer SHA-256: 1ee24339005b98f24afb2375d080f1aaa8ff1ab0e8812427b8b103b890cef98a
% Requer: pgfplots, biblioteca statistics (somente taxas.tex).
% Ajuste opcional no chamador: \def\AlturaGrafico{4.4cm} e \def\LarguraGrafico{\linewidth}.
\begingroup
\providecommand{\AlturaGrafico}{4.4cm}
\providecommand{\LarguraGrafico}{\linewidth}
\definecolor{grafverde}{HTML}{24574D}
\definecolor{grafazul}{HTML}{1769A5}
\definecolor{grafgrafite}{HTML}{253A40}
\definecolor{graflaranja}{HTML}{B97736}
\begin{tikzpicture}
\begin{axis}[width=\LarguraGrafico,height=\AlturaGrafico,
font=\scriptsize,tick label style={font=\scriptsize},
label style={font=\scriptsize},
axis line style={grafgrafite!65},tick style={grafgrafite!65},
grid style={black!10},
legend style={font=\scriptsize,draw=none,fill=none,cells={anchor=west}},
/pgf/number format/use comma,
xbar,bar width=4pt,xmin=0,xmax=158,
xtick={0,35,70,105,140},ytick={1,2,3,4,5},
yticklabels={Convers\~ao,Soja,Milho,Autos,Chuva},
y dir=reverse,ymin=0.4,ymax=5.6,xmajorgrids=true,
xlabel={Munic\'ipios},
legend columns=2,legend style={at={(0.5,-0.28)},anchor=north},
nodes near coords={\pgfmathprintnumber[fixed,precision=0]{\pgfplotspointmeta}},
nodes near coords align={horizontal},
every node near coord/.append style={font=\scriptsize,text=grafgrafite},
]
\addplot+[draw=grafverde,fill=grafverde] coordinates {(137,1) (136,2) (141,3) (139,4) (141,5)};
\addplot+[draw=grafazul,fill=grafazul] coordinates {(95,1) (66,2) (121,3) (4,4) (141,5)};
\legend{Algum dado,24 anos num\'ericos}
\end{axis}
\end{tikzpicture}
\endgroup
}
\column{.45\textwidth}
\small\textbf{Médias anuais municipais observadas}\par\smallskip
\begin{tabularx}{\linewidth}{L r r}
\toprule
\textbf{Medida} & \textbf{Média} & \textbf{Mediana}\\
\midrule
Autos (número) & 4,58 & 2,36\\
Multas (R\$ mi) & 6,08 & 1,40\\
\bottomrule
\end{tabularx}\par
\smallskip
\footnotesize N = 139; valores reais de dezembro de 2024.\par
\medskip\small
Entre os registros anuais numéricos, \textbf{60,4\%} dos autos e \textbf{60,6\%} das multas têm indicação de parcialidade.
\tpar
Valor autuado não equivale a pagamento; mais autos não demonstram maior efetividade.
\end{columns}
\end{frame}

% 15 — Parecer, seções 2, 4 e 7.
\begin{frame}{df2a: interpretação econômica dos agregados}
\begin{itemize}
\setlength{\itemsep}{4mm}
\item A predominância da pastagem na conversão direta justifica examinar os \textbf{incentivos específicos da pecuária}, mantendo agricultura e pastagem como destinos distintos.
\item A lotação aparente média municipal é de \textbf{1,471 cabeça/ha}. Ela descreve a relação entre rebanho e pastagem; não mede eficiência nem demonstra intensificação tecnológica.
\item A desigualdade e a parcialidade dos registros institucionais podem refletir infrações, detecção, direcionamento da fiscalização e registro. Os agregados \textbf{não separam esses componentes}.
\end{itemize}
\fonte{Os dados brutos não permitem inferir o agente, a legalidade ou o incentivo da conversão; fiscalização pode reagir à própria pressão ambiental.}
\end{frame}

% 16 — Parecer, seções 6 a 9.
\begin{frame}{Como df2a se relaciona com o artigo}
\small
\begin{tabularx}{\textwidth}{P{3.35cm} L}
\toprule
\textbf{Dimensão de H1} & \textbf{O que a análise preliminar indica para a investigação}\\
\midrule
Preços e capacidade & Entre os 95 municípios com Y completo, 39 têm capacidade histórica PAM de soja e 78 de milho. A cobertura conjunta restringe a comparação.\\[2mm]
Ambiente institucional & É necessário construir medidas válidas e distinguir condicionamento institucional de reação da fiscalização à conversão.\\[2mm]
Atividades e biomas & As diferenças entre destinos e territórios motivam comparações específicas, respeitando a localização das transições e a amostra disponível.\\
\bottomrule
\end{tabularx}\par
\medskip
\normalsize\textbf{Relação a avaliar:} como a conversão anual responde aos incentivos, conforme capacidade produtiva e condições institucionais.
\end{frame}

% 17 — Pendências explicitamente solicitadas pelo autor.
\begin{frame}{Próximas etapas}
\begin{block}{Delimitação territorial}
Ainda falta delimitar os municípios conforme \textbf{Cerrado, Amazônia, Pantanal e áreas ecótonas}, definindo como tratar territórios com mais de um bioma e localizar a conversão em cada contexto.
\end{block}
\medskip
\begin{block}{Pesquisa de campo em Mato Grosso}
Incluir os dados da pesquisa de campo junto à \textbf{SEMA}, ao \textbf{INTERMAT} e ao \textbf{INCRA}, aprofundando a caracterização das condições institucionais e territoriais.
\end{block}
\medskip
\centering\textcolor{verde}{Os resultados preliminares descrevem o fenômeno e orientam a continuidade da investigação.}
\end{frame}

\end{document}
